\newcommand\UniversiteAdi{Niğde Ömer Halisdemir Üniversitesi}
\newcommand\BolumAdi{MEKATRONİK BÖLÜMÜ}
\newcommand\DersKodu{MKT2002}
\newcommand\DersAdi{BİLGİSAYARLI KONTROL SİSTEMLERİ}
\newcommand\SinavAdi{Ara Sınav}
\newcommand\SinavTarihi{11.04.2025}
\newcommand\SinavSaati{14:00}
\newcommand\SinavSuresi{60dk}

\pagestyle{fancy}
\fancyhf{} % clear existing header/footer entries
\fancyhead[R]{Öğrenci No:\hspace{4.5cm}}
\fancyhead[L]{Ad Soyad:\hspace{7cm}}
\noindent \includegraphics[width=0.1\textwidth]{logo}
\begin{tabular}{
    p{0.15\linewidth}
    p{0.15\linewidth}
    p{0.2\linewidth}
    p{0.1\linewidth}
    p{0.15\linewidth}}
    \multicolumn{5}{c}{\textbf{\BolumAdi}}\\
    \multicolumn{5}{c}{\textbf{\DersAdi}}\\\hline
    \multicolumn{1}{|r|}{Ders Kodu:}&
    \multicolumn{1}{|c|}{\DersKodu}&
    \multicolumn{1}{|c|}{}& 
    \multicolumn{1}{|r|}{Tarih:}&
    \multicolumn{1}{|c|}{\SinavTarihi} \\\hline
    \multicolumn{1}{|r|}{Sınav Türü:}&
    \multicolumn{1}{|c|}{\SinavAdi}&  
    \multicolumn{1}{|c|}{}&
    \multicolumn{1}{|r|}{Saat:}&
    \multicolumn{1}{|c|}{\SinavSaati}\\\hline
    \multicolumn{1}{|r|}{Dönemi:}&
    \multicolumn{1}{|c|}{2024-2025}&
    \multicolumn{1}{|c|}{}&
    \multicolumn{1}{|r|}{Süre:}&
    \multicolumn{1}{|c|}{\SinavSuresi} \\\hline
    &&&&\\
\end{tabular}\\\\
\noindent\begin{center}
\begin{tabular}{|r|c|c|c|c|}\hline
    \textbf{Soru:}&
    \textbf{1}&
    \textbf{2}&
    \textbf{3}&
    \textbf{Toplam}\\\hline
    \textbf{Puan:}&
    \textbf{40}&
    \textbf{30}&
    \textbf{30}&
    \textbf{100}\\\hline
    \textbf{Not:}&&&&\\\hline
\end{tabular}\end{center}
\noindent\textbf{Uyarı:}
\begin{itemize}\bfseries
    \item Soruları dikkatlice okuyunuz. Hesap makinesi kullanılabilir.
    \item Defter, kitap ve notlar açık bir sınavdır.
    \item İşlemleri atlamadan ve ayrıntılı olarak veriniz. Sadece nümerik yanıtlar veya çizimler ara işlemler olmadan kabul edilmemektedir.
\end{itemize}

Bir transfer fonksiyonu
\begin{equation}
    G(s)=\frac{12}{s^2+5s+4}
\end{equation}
olarak verilmiştir.
\begin{enumerate}[\bfseries Soru-1]
    \item(40p)\, $G(s)$ için sönüm oranı $\zeta$'yı ve doğal frekans $\omega_n$'yi hesaplayınız.
    \item(30p)\, $G(s)$ sistemine birim basamak giriş
    \begin{equation}
        u(s)=\frac{1}{s}
    \end{equation}
    uygulanması durumunda oluşan çıkış $y(s)$'yi elde ediniz.
    \item(30p)\, Elde ettiğiniz $y(s)$'yi zaman tanım bölgesine çeviriniz. $y(t)=?$
\end{enumerate}